\documentclass[12pt,a4paper]{article}
\usepackage{amsmath}
\usepackage{amsfonts}
\usepackage{amssymb}
\author{Óscar Pozo de Cádiz}
\title{Especificación para el formato \texttt{Binary Scene Format (.bsf)}}

\usepackage{fancyhdr}
\pagestyle{fancy}
\lhead{Especificación \texttt{.bsf}}
\chead{Versión 0}
\rhead{\today}

\usepackage{tablefootnote}

%\usepackage[left=4cm,right=4cm,top=4cm,bottom=4cm]{geometry}

\begin{document}
\maketitle
\thispagestyle{fancy}

\section{Abstract}

\pagebreak

\section{Contenido}

El formato está organizado en bloques binarios, y sigue la siguiente estructura básica:

\subsection{Estructura general}

\begin{table}[h]
\center
\begin{tabular}{llll}
\hline
Offset & Tamaño & Tipo & Descripción \\ \hline
0 & 1 & \texttt{UINT8} & \textit{Magic number}. Debe ser \texttt{1}. \\
1 & 1 & \texttt{UINT8}& Versión de la especificación. Debe ser \texttt{0}. \\
2 & 4 & \texttt{UINT64} & Offset del bloque de sistemas. \\
6 & 4 & \texttt{UINT64} & Offset del bloque de objetos. \\
10 & 4 & \texttt{UINT64} & Offset del bloque de componentes. \\
14 & $s_{s}$\tablefootnote{Tamaño del bloque de sistemas.} & \texttt{Blob} & Bloque con la información de los sistemas. \\
14 + $s_{s}$ & $s_{o}$\tablefootnote{Tamaño del bloque de objetos.} & \texttt{Blob} & Bloque con la información de los objetos. \\
14 + $s_{s}$ + $s_{o}$& $s_{c}$\tablefootnote{Tamaño del bloque de componentes.} & \texttt{Blob} & Bloque con la información de los componentes. \\
\hline
\end{tabular}
\end{table}

\subsection{Bloque de sistemas}

\subsubsection{Estructura general}

\begin{table}[h]
\center
\begin{tabular}{llll}
\hline
Offset$_\text{bloque}$ \tablefootnote{Offset$_{bloque}$ = offset respecto al inicio del bloque, no del archivo.} & Tamaño & Tipo & Descripción \\ \hline
0 & 4 & \texttt{UINT64} & Número de sistemas. \\
4 & 4 & \texttt{UINT64} & Offset$_{bloque}$ del bloque de datos de sistemas. \\
8 & $s_{sn}$ & \texttt{STRING[]} & Nombres de sistemas (null-terminated). \\
8 + $s_{sn}$ & $s_{sd}$ & \texttt{Blob[]} & Bloques individuales de datos de sistemas. \\
\hline
\end{tabular}
\end{table}

\subsubsection{Bloque de datos de sistemas}

\begin{table}[h]
\center
\begin{tabular}{llll}
\hline
Offset$_\text{bloque}$& Tamaño & Tipo & Descripción \\ \hline
0 & 4 & \texttt{UINT64} & Tamaño del bloque. \\
4 & ... & \texttt{Blob} & Datos del sistema en cuestión. \\
\hline
\end{tabular}
\end{table}

\subsection{Bloque de objetos}

\begin{table}[h]
\center
\begin{tabular}{llll}
\hline
Offset$_\text{bloque}$& Tamaño & Tipo & Descripción \\ \hline
0 & 4 & \texttt{UINT64} & Número de objetos. \\
4 & ... & \texttt{UINT64[]} & Identificadores de los objetos. \\
\hline
\end{tabular}
\end{table}

\subsection{Bloque de componentes}

\begin{table}[h]
\center
\begin{tabular}{llll}
\hline
Offset$_\text{bloque}$ & Tamaño & Tipo & Descripción \\ \hline
0 & 4 & \texttt{UINT64} & Número de tipos de componentes. \\
4 & 4 & \texttt{UINT64} & Offset$_{bloque}$ del bloque de datos de componentes. \\
8 & $s_{cn}$\tablefootnote{Tamaño del bloque de nombres de tipos de componentes.} & \texttt{STRING[]} & Nombres de tipos de componentes (null-terminated). \\
12 + $s_{cn}$ & $s_{cd}$\tablefootnote{Tamaño del bloque de componentes.} & \texttt{Blob[]} & Bloques individuales de datos de componentes. \\
\hline
\end{tabular}
\end{table}

\subsection{Bloque individual de componentes}

\begin{table}[h]
\center
\begin{tabular}{llll}
\hline
Offset$_\text{bloque}$ & Tamaño & Tipo & Descripción \\ \hline
0 & 4 & \texttt{UINT64} & Número de componentes. \\
4 & 4 & \texttt{UINT64} & Offset$_{bloque}$ del bloque de IDs de objetos. \\
8 & $s_{cs}$\tablefootnote{Tamaño de todos los datos de componentes.} & \texttt{Blob[]} & Datos de componentes. \\
8 + $s_{cs}$ & $s_{o}$\tablefootnote{Tamaño de todos los identificadores de un tipo de componente.} & \texttt{UINT64[]} & Identificadores de objetos enlazados a los componentes. \\
\hline
\end{tabular}
\end{table}


\pagebreak
\appendix
\section{Bloques de sistemas por defecto}

\begin{table}[h]
\caption{OSK::ColliderRenderSystem}
\center
\begin{tabular}{llll}
\hline
Offset$_\text{bloque}$ & Tamaño & Tipo & Descripción \\ \hline
0 & 4 & \texttt{UINT64} & ID de la cámara. \\
\hline
\end{tabular}
\end{table}

\begin{table}[h]
\caption{OSK::RenderBoundsRenderer}
\center
\begin{tabular}{llll}
\hline
Offset$_\text{bloque}$ & Tamaño & Tipo & Descripción \\ \hline
0 & 4 & \texttt{UINT64} & ID de la cámara. \\
\hline
\end{tabular}
\end{table}

\begin{table}[h]
\caption{OSK::SkyboxRenderSystem}
\center
\begin{tabular}{llll}
\hline
Offset$_\text{bloque}$ & Tamaño & Tipo & Descripción \\ \hline
0 & 4 & \texttt{UINT64} & ID de la cámara. \\
\hline
\end{tabular}
\end{table}

\begin{table}[h]
\caption{OSK::TreeNormalsRenderSystem}
\center
\begin{tabular}{llll}
\hline
Offset$_\text{bloque}$ & Tamaño & Tipo & Descripción \\ \hline
0 & 4 & \texttt{UINT64} & ID de la cámara. \\
\hline
\end{tabular}
\end{table}

\begin{table}[h]
\caption{OSK::DeferredRenderSystem}
\center
\begin{tabular}{llll}
\hline
Offset$_\text{bloque}$ & Tamaño & Tipo & Descripción \\ \hline
0 & 4 & \texttt{UINT64} & ID de la cámara. \\
4 & 32 & \texttt{Blob} & Luz direccional (\texttt{DirectionalLigh}). \\
36 & 16 & \texttt{Blob} & Configuración IBL (\texttt{PbrIblConfig}). \\
52 & 52 & \texttt{Blob} & Mapa de sombras (\texttt{ShadowMap}). \\
104 & ... & \texttt{STRING} & Ruta del mapa IBL de irradiancia. \\
... & ... & \texttt{STRING} & Ruta del mapa IBL especular. \\
... & ... & \texttt{STRING} & Nombre del pase de resolución. \\
... & 4 & \texttt{UINT64} & Número de pases de renderizado. \\
... & ... & \texttt{STRING[]} & Nombres de los pases de renderizado. \\
... & 4 & \texttt{UINT64} & Número de pases de sombras. \\
... & ... & \texttt{STRING[]} & Nombres de los pases de sombras. \\
\hline
\end{tabular}
\end{table}

\begin{table}[h]
\caption{OSK::PhysicsSystem}
\center
\begin{tabular}{llll}
\hline
Offset$_\text{bloque}$ & Tamaño & Tipo & Descripción \\ \hline
0 & 12 & \texttt{Blob} & Aceleración de gravedad (\texttt{Vector3f}) (no usado). \\
12 & 12 & \texttt{Blob} & Dirección de gravedad (\texttt{Vector3f}) (no usado). \\
\hline
\end{tabular}
\end{table}

\pagebreak
\section{Bloques de componentes por defecto}

\begin{table}[h]
\caption{OSK::CameraComponent2D}
\center
\begin{tabular}{llll}
\hline
Offset$_\text{bloque}$ & Tamaño & Tipo & Descripción \\ \hline
0 & 64 & \texttt{Blob} & Matriz (\texttt{glm::mat4}). \\
\hline
\end{tabular}
\end{table}

\begin{table}[h]
\caption{OSK::CameraComponent3D}
\center
\begin{tabular}{llll}
\hline
Offset$_\text{bloque}$ & Tamaño & Tipo & Descripción \\ \hline
0 & 4 & \texttt{float} & Field of view. \\
4 & 4 & \texttt{float} & Límite inferior de field of view. \\
8 & 4 & \texttt{float} & Límite superior de field of view. \\
12 & 4 & \texttt{float} & Plano cercano. \\
16 & 4 & \texttt{float} & Plano lejano. \\
20 & 8 & \texttt{Blob} & Ángulos (\texttt{Vector2f}). \\
28 & 8 & \texttt{float} & Ángulos acumulados (\texttt{Vector2f}). \\
\hline
\end{tabular}
\end{table}

\begin{table}[h]
\caption{OSK::CollisionComponent}
\center
\begin{tabular}{llll}
\hline
Offset$_\text{bloque}$ & Tamaño & Tipo & Descripción \\ \hline
0 & 1 & \texttt{UINT8} & 1 si contiene colisionador cargado, 0 en caso contrario. \\
1 & 1 & \texttt{UINT8} & 1 si contiene colisionador customizado, 0 en caso contrario. \\
2 & ... & \texttt{STRING} & Ruta del archivo de colisión (null-terminated). \\
... & ... & \texttt{Blob} & Datos del colsionador customizado (\texttt{Collider}). \\
\hline
\end{tabular}
\end{table}

\begin{table}[h]
\caption{OSK::ModelComponent3D}
\center
\begin{tabular}{llll}
\hline
Offset$_\text{bloque}$ & Tamaño & Tipo & Descripción \\ \hline
0 & ... & \texttt{STRING} & Ruta del modelo 3Ds. \\
... & 1 & \texttt{UINT8} & 1 si proyecta sombras. \\
... & 4 & \texttt{UINT64} & Número de pases de shaders. \\
... & ... & \texttt{STRING[]} & Nombres de pases de shaders. \\
\hline
\end{tabular}
\end{table}

\begin{table}[h]
\caption{OSK::PhysicsComponent}
\center
\begin{tabular}{llll}
\hline
Offset$_\text{bloque}$ & Tamaño & Tipo & Descripción \\ \hline
0 & 1 & \texttt{UINT8} & 1 si el objeto es inmóvil. \\
1 & 4 & \texttt{float} & Masa si el objeto no es inmóvil, masa inversa en caso contrario. \\
5 & 36 & \texttt{Blob} & Tensor de inercia (\texttt{glm::mat3}). \\
41 & 12 & \texttt{Blob} & Velocidad. \\
53 & 12 & \texttt{Blob} & Aceleración. \\
65 & 12 & \texttt{Blob} & Velocidad angular. \\
\hline
\end{tabular}
\end{table}

\begin{table}[h]
\caption{OSK::Transform2D}
\center
\begin{tabular}{llll}
\hline
Offset$_\text{bloque}$ & Tamaño & Tipo & Descripción \\ \hline
0 & 4 & \texttt{UINT64} & ID del objeto que posee el componente. \\
4 & 8 & \texttt{UINT64} & ID del objeto padre. \\
12 & 8 & \texttt{Blob} & Posición (\texttt{Vector2f}). \\
20 & 8 & \texttt{Blob} & Escala (\texttt{Vector2f}). \\
28 & 4 & \texttt{Blob} & Rotación. \\
32 & 4 & \texttt{UINT64} & Número de objetos hijos. \\
36 & ... & \texttt{UINT64[]} & ID de los objetos hijos. \\
\hline
\end{tabular}
\end{table}

\begin{table}[h]
\caption{OSK::Transform3D}
\center
\begin{tabular}{llll}
\hline
Offset$_\text{bloque}$ & Tamaño & Tipo & Descripción \\ \hline
0 & 4 & \texttt{UINT64} & ID del objeto que posee el componente. \\
4 & 1 & \texttt{UINT8} & 1 si hereda la posición. \\
5 & 1 & \texttt{UINT8} & 1 si hereda la rotación. \\
6 & 1 & \texttt{UINT8} & 1 si hereda la escala. \\
7 & 12 & \texttt{Blob} & Posición local (\texttt{Vector3f}). \\
19 & 64 & \texttt{Blob} & Rotación local (\texttt{Quaterion}). \\
83 & 12 & \texttt{Blob} & Escala local (\texttt{Vector3f}). \\
95 & 4 & \texttt{UINT64} & Número de elementos hijos. \\
99 & ... & \texttt{UINT64[]} & ID de los objetos hijos. \\
\hline
\end{tabular}
\end{table}


\pagebreak
\section{Bloques de datos por defecto}

\begin{table}[h]
\caption{OSK::GRAPHICS::PbrIblConfig}
\center
\begin{tabular}{llll}
\hline
Offset$_\text{bloque}$ & Tamaño & Tipo & Descripción \\ \hline
0 & 4 & \texttt{float} & Fuerza de mapa de irradiancia. \\
4 & 4 & \texttt{float} & Fuerza de mapa especular. \\
8 & 4 & \texttt{float} & Fuerza de radiancia. \\
12 & 4 & \texttt{float} & Fuerza de materiales emisivos. \\
\hline
\end{tabular}
\end{table}

\begin{table}[h]
\caption{OSK::GRAPHICS::DirectionalLight}
\center
\begin{tabular}{llll}
\hline
Offset$_\text{bloque}$ & Tamaño & Tipo & Descripción \\ \hline
0 & 16 & \texttt{Blob} & Color de la luz (\texttt{Color}). \\
16 & 12 & \texttt{Blob} & Dirección de la luz (\texttt{Vector3f}). \\
28 & 4 & \texttt{float} & Fuerza de la luz. \\
\hline
\end{tabular}
\end{table}

\begin{table}[h]
\caption{OSK::GRAPHICS::ShadowMap}
\center
\begin{tabular}{llll}
\hline
Offset$_\text{bloque}$ & Tamaño & Tipo & Descripción \\ \hline
0 & 4 & \texttt{UINT64} & Identificador de la cámara. \\
4 & 4 & \texttt{float} & Distancia del primer split. \\
8 & 4 & \texttt{float} & Distancia del segundo split. \\
12 & 4 & \texttt{float} & Distancia del tercer split. \\
16 & 4 & \texttt{float} & Distancia del cuarto split. \\
20 & 4 & \texttt{float} & Distancia hasta el near plane. \\
24 & 4 & \texttt{float} & Distancia hasta el far plane. \\
28 & 12 & \texttt{Blob} & Dirección de la luz (\texttt{Vector3f}). \\
40 & 12 & \texttt{Blob} & Punto de la luz (\texttt{Vector3f}). \\
\hline
\end{tabular}
\end{table}

\end{document}